% Work-MAT | 2. rocnik | pripravovana hodina 31
% Návaznost: uživatelské připravy.zip (hodiny 01–20), autorská učebnice na GitHubu.
% Realisticky.cz: didaktická inspirace v posloupnosti a typech úloh; zadání jsou vlastní.
\documentclass[11pt,a4paper]{article}
\input{../../sablony/pripravy/priprava.sty}
\nastavPripravu{2. ročník}{Exponenciální rovnice}{31}{Souhrnné opakování}
\begin{document}
\blok{Motivace a návaznost}
Poslední hodina před logaritmy propojí všechny dosud známé metody; žáci před každou úlohou pojmenují postup. Začátek věnuj sebereflexi po práci P1.
\poznamka{Orientační nabídka pro 45--55 minut: krátké zopakování a motivace, nový výklad, tři společné vzorové úlohy, samostatná práce. Vyber podle tempa třídy; zbývající zadání slouží jako rezerva.}
\blok{Výklad}
Vytýkání rozpoznáme podle součtu posunutých mocnin. Substituce se hodí pro opakující se mocninu $a^x$ a její vyšší mocniny; nerovnice řídí monotónnost.
\blok{Společné příklady}
\begin{prikladReseniBox}[title={Příklad 1}]
Vyřeš, resp. urči: $2^{x+1}+2^x=24$.
\tcblower
\[\begin{aligned}3\cdot2^x=24 \\ x=3\end{aligned}\]
\end{prikladReseniBox}
\begin{prikladReseniBox}[title={Příklad 2}]
Vyřeš, resp. urči: $4^x-5\cdot2^x+4=0$.
\tcblower
\[\begin{aligned}t=2^x>0,\quad (t-1)(t-4)=0 \\ K=\{0,2\}\end{aligned}\]
\end{prikladReseniBox}
\begin{prikladReseniBox}[title={Příklad 3}]
Vyřeš, resp. urči: $(\frac{1}{2})^x<(\frac{1}{2})^3$.
\tcblower
\[\begin{aligned}x>3\end{aligned}\]
\end{prikladReseniBox}
\poznamka{Zpětná vazba může zabrat 10 minut, zbývající příklady jsou rezerva; žákům nezadávej všechny mechanicky.}
\blok{Samostatné procvičování}
\begin{zadaniReseniBox}[title={Samostatná práce 1}]
\begin{enumerate}[label=\alph*)]
\item $3^x+3^{x+1}=108$.
\item $9^x-4\cdot3^x+3=0$.
\item $2^x+2^{3-x}=6$.
\item $8^{x-1}\le4^{x+1}$.
\item $3^x+3^x=\frac{2}{3}$.
\item $4^x-10\cdot2^x+16=0$.
\end{enumerate}
\tcblower
\begin{enumerate}[label=\alph*)]
\item $x=3$.
\item $K=\{0,1\}$.
\item $K=\{1,2\}$.
\item $x\le5$.
\item $x=-1$.
\item $K=\{1,3\}$.
\end{enumerate}
\end{zadaniReseniBox}
\blok{Dobrovolné rozšíření}
\blok{Úkol na doma}
Při početních úlohách napiš postup, zvolenou metodu a zkoušku v původním zadání. U soustav kontroluj obě rovnice; u nerovnic výsledek zakresli na osu.
\begin{zadaniReseniBox}[title={Domácí úkol}]
\begin{enumerate}[label=\alph*)]
\item $2^{x+2}-2^x=12$.
\item $9^x-10\cdot3^x+9=0$.
\item $2^x+2^{3-x}=6$.
\item $3^{x+1}+3^x=108$.
\item $8^{x-1}\le4^{x+1}$.
\item $4^x-5\cdot2^x+4=0$.
\item Kdy vyřadíme pomocný kořen?.
\end{enumerate}
\tcblower
\begin{enumerate}[label=\alph*)]
\item $x=2$.
\item $K=\{0,2\}$.
\item $K=\{1,2\}$.
\item $x=3$.
\item $x\le5$.
\item $K=\{0,2\}$.
\item Pokud $t\le0$..
\end{enumerate}
\end{zadaniReseniBox}
\poznamka{Domácí úkol obsahuje 7 položek; předpokládaná domácí práce přibližně 15--30 minut. Po větších blocích jsou vědomě promíchané dříve probrané typy.}
\end{document}
